
\documentclass{article}
\usepackage{colortbl}
\usepackage{makecell}
\usepackage{multirow}
\usepackage{supertabular}

\begin{document}

\newcounter{utterance}

\centering \large Interaction Transcript for game `ifeval', experiment `ifeval', episode 54 with reflexion{-}Qwen3.5{-}4B.
\vspace{24pt}

{ \footnotesize  \setcounter{utterance}{1}
\setlength{\tabcolsep}{0pt}
\begin{supertabular}{c@{$\;$}|p{.15\linewidth}@{}p{.15\linewidth}p{.15\linewidth}p{.15\linewidth}p{.15\linewidth}p{.15\linewidth}}
    \# & \multicolumn{2}{c}{Player} && \multicolumn{2}{c}{Game Master} \\
    \hline

    \theutterance \stepcounter{utterance}  
    & & & \multicolumn{4}{p{0.6\linewidth}}{
        \cellcolor[rgb]{0.9,0.9,0.9}{
            \makecell[{{p{\linewidth}}}]{
                \texttt{\tiny{[P1$\langle$GM]}}
                \texttt{Take the text below as a starting point, and make it a complete article: "You may have to meet with a helper to work out a parenting plan. The first would be to estimate how many time you have everyday for parenting, and is that enough...."} \\
\texttt{Avoid using the following keywords: sleep, cook, feed} \\
\texttt{Mention the keyword 'schedule' for more than 5 times.} \\
            }
        }
    }
     \\ \\

    \theutterance \stepcounter{utterance}  
    & \multicolumn{4}{p{0.6\linewidth}}{
        \cellcolor[rgb]{0.561,0.737,0.561}{
            \makecell[{{p{\linewidth}}}]{
                \texttt{\tiny{[P1$\rangle$GM]}}
                \texttt{\# Navigating the Complexity of Co{-}Parenting: Creating a Functional Plan} \\
\\ 
\texttt{You may have to meet with a helper to work out a parenting plan. The first step in this process is to estimate how many hours you have everyday for parenting, and is that enough. This initial assessment sets the foundation for a sustainable arrangement that prioritizes the well{-}being of the child while respecting the logistical realities of both parents' lives. Without a clear understanding of available time, any agreement made is likely to fail under the pressure of daily demands.} \\
\\ 
\texttt{Once the time estimate is established, the next critical phase involves drafting a detailed **schedule**. This **schedule** must account for school hours, extracurricular activities, and essential appointments. A flexible **schedule** allows for adjustments when unexpected events arise, ensuring that neither parent feels overwhelmed by rigid commitments. However, a **schedule** is only as good as its execution; it requires consistent communication and mutual respect to remain effective.} \\
\\ 
\texttt{Consider the specific needs of the child when finalizing the **schedule**. Does the current **schedule** provide adequate time for homework, play, and rest? If the answer is no, the **schedule** must be revised to better balance work and family life. Many families struggle because they overlook the importance of a realistic **schedule** that accommodates the child's developmental stage. A well{-}structured **schedule** not only reduces conflict between parents but also provides the child with a sense of stability and predictability.} \\
\\ 
\texttt{Ultimately, the goal is to create a **schedule** that fosters a healthy environment for growth and learning. Whether you are working with a mediator or negotiating directly with a co{-}parent, the focus should remain on the child's best interests. By carefully estimating time and adhering to a thoughtful **schedule**, parents can navigate the challenges of co{-}parenting with confidence and clarity. Remember, a successful **schedule** is one that evolves with the child's changing needs, ensuring long{-}term harmony in the family unit.} \\
            }
        }
    }
    & & \\ \\

    \theutterance \stepcounter{utterance}  
    & & & \multicolumn{2}{p{0.3\linewidth}}{
        \cellcolor[rgb]{0.9,0.9,0.9}{
            \makecell[{{p{\linewidth}}}]{
                \texttt{\tiny{[GM$|$GM]}}
                {'keywords:forbidden_words': {'forbidden_words': ['sleep', 'cook', 'feed']}, 'keywords:frequency': {'relation': 'at least', 'keyword': 'schedule', 'frequency': 6}}
            }
        }
    }
    & & \\ \\

    \theutterance \stepcounter{utterance}  
    & & & \multicolumn{2}{p{0.3\linewidth}}{
        \cellcolor[rgb]{0.9,0.9,0.9}{
            \makecell[{{p{\linewidth}}}]{
                \texttt{\tiny{[GM$|$GM]}}
                \texttt{game\_result = WIN} \\
            }
        }
    }
    & & \\ \\

\end{supertabular}
}

\end{document}
